\documentclass[10pt,a4paper]{amsart}
\usepackage[margin=19mm]{geometry}
\usepackage{amsmath,amssymb,mathtools}
\usepackage[scr=rsfs,cal=euler]{mathalfa}
\usepackage{xfrac}
\usepackage[unicode,hidelinks,bookmarks=false]{hyperref}
\newcommand{\ie}{i.e.\ }
\newcommand{\cf}{cf.\ }
\newcommand{\resp}{respectively\ }
\newcommand{\Subsection}{\subsection}
\providecommand{\nobreakdash}{\nobreak\hbox{-}}
\setlength{\emergencystretch}{2em}
\multlinegap0pt
\usepackage{bm}
\usepackage{bbm}
\usepackage{dsfont}
\usepackage{xspace}
\usepackage{cancel}
\usepackage{mathtools}
\usepackage{amsthm}
\makeatletter
\@ifundefined{definition}{\newtheorem{definition}{Definition}}{}
\@ifundefined{proposition}{\newtheorem{proposition}[definition]{Proposition}}{}
\makeatother
\newcommand{\defeq}{\vcentcolon=}
\title[Nonlinear Schr\"odinger Equation with magnetic potential on ]{Nonlinear Schr\"odinger Equation with magnetic potential on metric graphs}
\author{Nicol\`o Cangiotti, Ivan Gallo, David Spitzkopf}
\date{Source: arXiv:2601.23115v2 (CC BY 4.0)}
\begin{document}
\maketitle

\noindent\emph{Source and licence.} Nonlinear Schr\"odinger Equation with magnetic potential on metric graphs. Version arXiv:2601.23115v2 (CC BY 4.0), distributed under the Creative Commons Attribution 4.0 International licence, as recorded in the arXiv licence metadata for this exact version and shown on the abstract page https://arxiv.org/abs/2601.23115v2. Full rights notice: licenses/arxiv-2601.23115-CC-BY-4.0.txt.\par\smallskip

\noindent\textit{Editorial scope.} Counted unit: the Proposition whose source label is \texttt{equivalence} in the article (source0.tex lines 326--341, source label \texttt{equivalence}), delivered together with the article's own complete original proof of it (source0.tex lines 342--417, one \texttt{proof} environment of 76 lines). No mathematics is written, repaired or re-derived: statement and proof are the article's own text line for line. Required context retained uncounted: bc (lines 170--176). Every definition, display and label of those lines is retained unchanged. Omitted from this package and not redistributed: 1--169, 177--325, 418--955. Nothing inside a retained excerpt is omitted or altered: every retained body is delivered exactly as the article's lines are written, so no omission receipt of any kind accompanies this package. Citations: the retained excerpts carry no citation command at all, so no bibliography entry of the article is transcribed and no reference is redistributed. Reference adaptation: none was needed and none was made; every reference inside the delivered unit resolves inside it, so no unresolved reference or citation is produced. Printed slips kept literal: no printed word, symbol, hypothesis or number of the counted statement or of its proof was altered. Layout and class: the article's own class is \texttt{amsart} with packages including amsthm, graphicx, tikz, caption, subcaption, lipsum, geometry, xcolor, adjustbox, mathtools, leftindex, comment, bm, lineno; the wrapper is the lane's pinned \texttt{amsart} editorial wrapper at 10pt on A4, which restates the theorem-like environments the retained text needs and adds the article's own macro definitions that the retained text needs (\texttt{\textbackslash defeq}), with extra wrapper packages (bm, bbm, dsfont, xspace, cancel, mathtools). The delivered numbering is the unit's own. Assets: no retained span contains a figure environment, a graphics-inclusion command, a PDF-page inclusion, a table or a TikZ picture; no image file is copied into this package, the package declares no asset dependency and no image file is redistributed with it. Licence: version arXiv:2601.23115v2 of this article is distributed under the Creative Commons Attribution 4.0 International licence, as recorded in the arXiv licence metadata for this exact version and shown on the abstract page https://arxiv.org/abs/2601.23115v2. A full rights notice is delivered at licenses/arxiv-2601.23115-CC-BY-4.0.txt. Characters: every retained excerpt is pure ASCII, so the delivered engine, XeLaTeX with its default Latin Modern text font, reports no missing character.
\begin{equation}
    \label{bc}
    \begin{cases}
        u(x) \ \text{is continuous at every vertex } v \in \mathbb{V}_\mathcal{G}\,; \\
        \sum_{e \succ v} \frac{\textrm{d} u_e}{\textrm{d}x_e}(v) - i A_e(v) u(v) = 0 \quad \forall v \in \mathbb{V}_\mathcal{G}\,,
    \end{cases}   
\end{equation}

\begin{proposition} \label{equivalence}
Let $A \in C^1(\mathcal{G}) \cap L^2_{loc}(\mathcal{G})$. Then the minimization problem
    \begin{equation}
         \mathcal{E}_{A,\mathcal{G}}(\mu) \defeq \inf_{u \in H^1_{A,\mu}(\mathcal{G})} E_A(u, \mathcal{G})
    \end{equation}
    is equivalent to the problem
    \begin{equation}
        \mathcal{I}_\mathcal{G}(\mu) \defeq \inf_{v \in H^1_{\mu}(\mathcal{G})} I_A(v, \mathcal{G})\,,
        \label{Ig2}
    \end{equation}
    where the effective functional $I_A$ is defined as
    \begin{equation} \label{Ia}
        I_A(v,\mathcal{G})  = \frac{1}{2}\int_{\mathcal{G}} |v'|^2 \mathrm{d}x + \int_{C(\mathcal{G})} \Phi_{\gamma}(A) |v|^2 \mathrm{d}x - \frac{1}{p} \int_{\mathcal{G}} |v|^p \mathrm{d}x\,,
    \end{equation}
    with $C(\mathcal{G})$ denoting the set of independent cycles and $\Phi_{\gamma}(A)$ being a scalar repulsive potential determined by the magnetic flux.
\end{proposition}

\begin{proof}
Consider a function $u \in H^1_A(\mathcal{G})$. It can be decomposed as $u(x) = v(x)e^{i \theta(x)}$ with $v = |u|$.
A direct computation yields
\begin{align*}
    \int_{\mathcal{G}} |D u|^2 &= \int_\mathcal{G} \left| v'e^{i \theta} + i \theta' v e^{i\theta} - i A v e^{i\theta}\right|^2 \mathrm{d}x \\
    &= \int_\mathcal{G} \left| v' - i(A-\theta')v \right|^2 \mathrm{d}x \\
    &= \int_{\mathcal{G}} |v'|^2 \mathrm{d}x + \int_{\mathcal{G}} (A - \theta')^2 |v|^2 \mathrm{d}x\,.
\end{align*}
Since our goal is to characterize the solutions of $\mathcal{E}_{\mathcal{G}}(\mu)$, we reformulate the problem as
\begin{equation}
    \inf_{v \in H^1_{\mu}(\mathcal{G}), \hspace{1mm} \theta \in C^1(\mathcal{G}) } E_A(u, \mathcal{G}) \,,
\end{equation}
thereby decoupling the variables and focusing on the minimization with respect to $\theta$. In particular, we seek the function that minimizes the quantity
$$
\int_\mathcal{G} (A - \theta')^2 |v|^2 \mathrm{d}x\,,
$$
while preserving the boundary conditions \eqref{bc}. 
\bigskip

Let $\gamma$ be a cycle parameterized by $[0,L]$. While the continuity conditions at the vertices require $u(0) = u(L)$, the phase function $\theta(x)$ must only satisfy the single-valued condition:
\begin{equation}
\theta(L) - \theta(0) = 2\pi m, \hspace{0.5cm} m \in \mathbb{Z}\,.
\end{equation}

Using the \emph{Lagrange Multiplier Theorem} for fixed $v$, one can find the optimal $\theta'$ subject to the periodicity constraint 
\begin{equation}
    \int_0^L \theta'(x) \mathrm{d}x = 2 \pi m.
\end{equation}
For some $\lambda \in \mathbb{R}$, we have the variation:
\begin{equation}
\nabla (A-\theta')^2 - \lambda \nabla\left(\int_0^L \theta'(x) \mathrm{d}x - 2\pi m\right ) = 0\,.
\end{equation}
Therefore we set $\theta'(x) = A(x) + \frac{1}{2}\lambda L$ for some constant $\lambda$ to be determined. The single-valued constraint becomes:
\begin{equation}
\int_0^L \theta'(x)\, \mathrm{d}x = \int_0^L \left[A(x) + \frac{1}{2}\lambda L\right ]\, \mathrm{d}x = 2\pi m\,.
\end{equation}
This gives us:
\begin{equation}
\int_0^L A(x) \mathrm{d}x + \frac{1}{2}\lambda L^2 = 2\pi m\,,
\end{equation}
which yields:
\begin{equation}
\lambda = \frac{4\pi m}{L^2} - \frac{2}{L^2}\int_0^L A(x) \mathrm{d}x \,.
\end{equation}
Consequently, substituting $\lambda$ back into the expression for $(A-\theta')$, we obtain:
\begin{equation}
     A(x) - \theta'(x) = A(x) - \left(A(x) + \frac{1}{2}\lambda L\right) = \frac{1}{L}\int_0^L A(x) \mathrm{d}x - \frac{2\pi m}{L}\,,
\end{equation} 
and the magnetic energy contribution becomes:
\begin{equation}
\int_0^L |A - \theta'|^2|v|^2 \mathrm{d}x =  \int_0^L  \left(\frac{1}{L}\int_0^L A(x) \mathrm{d}x - \frac{2\pi m}{L}\right)^2|v|^2 \mathrm{d}x\,.
\end{equation}
Finally, we define $m$ as the integer satisfying the minimization condition:
$$
(\alpha_\gamma - m)^2 = \mathrm{dist}\left(\alpha_\gamma, \mathbb{Z}\right)^2 \defeq \min_{k \in \mathbb{Z}} (\alpha_\gamma - k)^2,
$$
where we have set $\alpha_\gamma \defeq \frac{1}{2\pi} \int_{0}^L A(x) \mathrm{d} x$.
And thus we have:
\begin{equation}
    \int_0^L \left|A - \theta'\right|^2|v|^2 \mathrm{d}x = \int_0^L \frac{4 \pi^2}{L^2} \mathrm{dist}\left( \frac{\int_0^L A \mathrm{d}x}{2 \pi}, \mathbb{Z} \right)^2 |v|^2 \mathrm{d}x\, .
\end{equation}
Then, $\left(A(x) - \theta'(x)\right)$ is constant on the loop, and depends only on the size of the loop and the average magnetic flux through the cycle. 

\medskip

We can now easily generalize this analysis for a general graph with more cycles. 
Let $C(\mathcal{G})$ be the set of independent cycles of $\mathcal{G}$. Then for every $\gamma \in C(\mathcal{G})$:
\begin{equation}
    \Phi_{\gamma}(A) = \frac{4 \pi^2}{|\gamma|^2} \mathrm{dist}\left( \frac{\int_\gamma A(x) \mathrm{d}x}{2 \pi}, \mathbb{Z} \right)^2
\end{equation}
and the energy functional becomes:
\begin{equation} \label{Ia2}
     I_A(v,\mathcal{G})  = \frac{1}{2}\int_{\mathcal{G}} |v'|^2 \mathrm{d}x + \int_{C(\mathcal{G})}\Phi_{\gamma}(A) |v|^2 \mathrm{d}x - \frac{1}{p} \int_{\mathcal{G}} |v|^p \mathrm{d}x\,.
\end{equation}
Since in the definition of Eq. \eqref{Ia} only $v = |u|$ appears, for Proposition \ref{equivalence} we have $v \in H^1_{\mu}(\mathcal{G})$ and the result follows.
\end{proof}

\end{document}
